\noindent
Brains are wired in patterns that look far from random, and a widely repeated claim holds
that those patterns are what make brains good at computing. How tightly a network's
wiring is tied to what it can compute is what we call its computational coupling, and it
cannot be read off a single network: it only shows up as a difference. The usual way to
test for it is to take a real wiring diagram, shuffle its connections at random, and check
whether the real one performs better at some task. If it does, the wiring pattern is
credited with the difference.

We tried to check one such result and found that the test has a loose part. ``Shuffle the
connections at random'' sounds like one procedure, but it is at least two, because a
connection has both a location and a strength, and the phrase does not say whether a
connection's strength should travel with it when it moves. Both readings are reasonable.
We wrote one of them, believing we had written the other.

On the same wiring diagrams, measured the same way, the two shuffling procedures disagree
about whether real brain wiring wins --- not by a little, but in opposite directions, each
with the statistical confidence normally taken as decisive. We can say why: one of the two
procedures quietly creates a pattern that real brains do not have, which handicaps the
comparison. That explanation makes a prediction we could have failed --- on a wiring
diagram where all connections have equal strength, there is nothing to move, so the two
procedures must agree. One of our five networks is of that kind, and there they agree
almost exactly.

Three further choices in the same test have the same character: what kind of model neuron
is used, how strongly the network is driven, and which neurons are allowed to receive
input and be read out. Each can reverse the answer on identical data, and none is usually
reported as a choice at all.

None of this shows that brain wiring is unremarkable. It shows that computational
coupling is not something a wiring diagram simply has --- it is something a wiring diagram
has relative to a way of asking, and a popular way of asking is less settled than it
appears, and that the code
which generates the comparison matters as much as the sentence describing it. We include
our own mistake in full, because it produced a clean and convincing result that was
answering the wrong question, and nothing in the output said so.
